\begin{abstract}
\noindent
We present \textbf{NepGraph}, an open-source analytical framework for
discovering organised market structure in the Nepal Stock Exchange (NEPSE)
from historical price data alone.  The pipeline transforms daily closing
prices into log returns, builds a Pearson correlation matrix over a
124-instrument universe spanning June 2003 to April 2026 (almost
23~years of NEPSE history), converts that matrix into a distance metric
$d_{ij}=\sqrt{2(1-\rho_{ij})}$, sparsifies via a top-$k$
nearest-neighbour graph, and applies the Louvain modularity-maximisation
algorithm on the resulting graph to identify hidden communities.  A
central methodological contribution is a pre-registered $k$-sweep over
\num{440} parameter combinations (eleven values of $k\in\{2,\dots,20\}$,
four temporal windows, ten Louvain seeds) that empirically justifies
$k=5$ as the \emph{minimum} sparsification that keeps the entire NEPSE
universe connected as a single graph.  While smaller sparsifications
(notably $k=2$ and $k=3$) yield a slightly higher headline modularity
score, they fragment the graph into multiple disconnected
\emph{islands} of stocks, which breaks down global topology metrics such
as betweenness centrality: one cannot measure a stock's ability to act
as a bridge between communities if there is, by construction, no path
between those communities in the first place.  Choosing $k=5$ therefore
represents the exact mathematical \emph{sweet spot}---the minimum
threshold required to keep the market connected as a single graph while
still maintaining high modularity ($Q=0.66$) and avoiding the
over-connected \emph{hairball} that larger $k$ values would produce.
The detected communities overlap only partially with NEPSE's thirteen
official sector labels (Adjusted Rand Index $\text{ARI}=0.21$); we
exploit this divergence as a \emph{sector-mismatch anomaly score} that
surfaces stocks whose dominant neighbours come from a different official
sector---candidates for hidden common-factor exposure.  Statistical
significance is established by an edge-swap null model ($z=47.0$,
$p<10^{-3}$), temporal stability by rolling 1-year ARI/NMI, and economic
interpretability by a four-quadrant centrality analysis (Spearman
$\rho \approx 0.08$ between eigenvector and betweenness, indicating the two
measures are essentially orthogonal).  A rolling out-of-sample backtest
(50~bps transaction costs) compares the detected-community baskets
against sector-based and equal-weight strategies.  All code, raw
metrics, and figures are released under MIT.
\end{abstract}
